\chapter{编译期计算、type traits 和约束边界}

\section{问题入口：有些错误能不能提前到编译期}

运行期检查很重要，但有些规则在编译期就知道。例如固定容量缓冲区的容量不能为 \code{0}。如果等运行到取模时才失败，错误出现得太晚。模板参数和 \code{constexpr} 让一部分规则可以提前：

\begin{lstlisting}[language=C++,caption={编译期容量检查}]
template <typename T, std::size_t CAPACITY>
    requires(CAPACITY > 0)
class RingBuffer {
    // ...
};
\end{lstlisting}

这不是为了炫技，而是把“不可能正确”的类型直接挡在编译期。编译期计算的价值在于更早暴露错误、更清楚表达常量关系，以及减少运行期分支。

\section{constexpr 函数}

\code{constexpr} 函数可以在编译期求值，也可以在运行期调用：

\begin{lstlisting}[language=C++,caption={端口范围检查}]
constexpr bool is_valid_port(int port) noexcept
{
    return 1 <= port && port <= 65535;
}

static_assert(is_valid_port(80));
static_assert(!is_valid_port(70000));
\end{lstlisting}

\code{static_assert} 是编译期断言。它适合检查常量规则、模板约束和平台假设。不要把运行期输入塞进 \code{static_assert}，因为运行期输入编译时并不存在。

\section{用 type traits 观察类型能力}

type traits 是一组在编译期回答类型问题的工具。例如：

\begin{lstlisting}[language=C++,caption={检查类型属性}]
static_assert(std::is_integral_v<int>);
static_assert(!std::is_integral_v<std::string>);
static_assert(std::is_move_constructible_v<std::vector<int>>);
\end{lstlisting}

这些工具常用于泛型库内部。普通业务代码不应该到处堆 traits；如果可以用 concept 表达接口，优先用 concept。traits 更像底层零件，concept 更像公开合同。

\section{enable\_if 到 concept 的演进}

\cpp20 以前常用 \code{std::enable_if} 限制模板：

\begin{lstlisting}[language=C++,caption={传统 SFINAE 写法}]
template <typename T, typename = std::enable_if_t<std::is_integral_v<T>>>
bool is_even(T value)
{
    return value % 2 == 0;
}
\end{lstlisting}

\cpp20 可以写得更直接：

\begin{lstlisting}[language=C++,caption={concept 写法更接近意图}]
template <std::integral T>
bool is_even(T value)
{
    return value % 2 == 0;
}
\end{lstlisting}

两者本质都在限制类型，但 concept 把错误消息和接口意图都变清楚。学习 traits 是为了能读懂旧代码和写底层库，不是为了把所有新代码写复杂。

\section{if constexpr：编译期分支}

有些分支取决于类型，不取决于运行期值：

\begin{lstlisting}[language=C++,caption={根据类型选择实现}]
template <typename T>
std::string to_debug_string(const T& value)
{
    if constexpr (std::is_same_v<T, std::string>) {
        return value;
    } else if constexpr (std::is_arithmetic_v<T>) {
        return std::to_string(value);
    } else {
        return "<object>";
    }
}
\end{lstlisting}

\code{if constexpr} 未选择的分支不会被实例化。这和普通 \code{if} 不同。普通 \code{if} 的两个分支都必须语法和类型正确；\code{if constexpr} 可以根据类型裁掉不适用分支。

\section{编译期表和运行期表}

如果一组数据固定不变，可以用 \code{constexpr std::array}：

\begin{lstlisting}[language=C++,caption={编译期关键字表}]
constexpr std::array<std::string_view, 4> KEYWORDS{
    "class",
    "const",
    "return",
    "while",
};

constexpr bool is_keyword(std::string_view word)
{
    for (std::string_view keyword : KEYWORDS) {
        if (keyword == word) {
            return true;
        }
    }
    return false;
}
\end{lstlisting}

这里的表很小，线性扫描足够。如果关键字很多，可以排序后二分，但不要为了“编译期”牺牲可读性。编译期优化也要有规模依据。

\section{模板错误如何设计得更友好}

一个泛型函数如果没有约束，用户传错类型时可能看到几十行错误。用 concept 或 \code{static_assert} 可以把错误变近：

\begin{lstlisting}[language=C++,caption={用编译期断言给出明确错误}]
template <typename T>
T median(std::span<const T> values)
{
    static_assert(std::totally_ordered<T>,
                  "median requires an ordered value type");
    // ...
}
\end{lstlisting}

更好的方式是直接写到模板头上：

\begin{lstlisting}[language=C++,caption={约束排序能力}]
template <std::totally_ordered T>
T median(std::span<const T> values);
\end{lstlisting}

公开接口优先用约束，函数内部再用 \code{static_assert} 检查更细的不变量。

\section{不要把运行期问题硬搬到编译期}

编译期计算有成本：模板实例化更复杂，编译时间可能增加，错误信息可能变长。配置文件、用户输入、网络数据、数据库记录都天然是运行期信息，不要为了“高级”硬做成模板参数。

一个实用判断：如果错误和类型或固定常量有关，考虑编译期；如果错误和外部输入有关，运行期处理更自然。

\begin{exercisebox}
实现一个 \code{FixedString<N>}，保存固定长度字符串，并写一个 \code{constexpr} 函数判断它是否为空。然后回答：哪些检查发生在编译期，哪些仍然只能运行期处理？
\end{exercisebox}

\section{constexpr 的普通代码理解}

\code{constexpr} 函数不是只能在编译期运行的函数。它更像一段同时满足“可在编译期求值条件”的普通函数：

\begin{lstlisting}[language=C++,caption={constexpr 可以编译期也可以运行期}]
constexpr bool is_valid_port(int port) noexcept
{
    return 1 <= port && port <= 65535;
}

static_assert(is_valid_port(80));      // 编译期
bool ok = is_valid_port(user_input);   // 运行期
\end{lstlisting}

当实参是常量表达式时，编译器可以在编译期计算；当实参来自用户输入时，只能运行期计算。不要把 \code{constexpr} 理解成“强制提前执行”。它提供的是提前执行的可能性。

\section{static\_assert 背后是编译期 if}

\code{static_assert} 可以理解成编译期门禁：

\begin{lstlisting}[language=C++,caption={平台假设检查}]
static_assert(sizeof(void*) == 8, "this example expects 64-bit pointers");
\end{lstlisting}

如果条件为假，编译停止，程序不会生成。它适合检查类型、模板参数、平台假设和固定常量。不适合检查配置文件端口，因为配置文件内容编译期并不存在。

\section{if constexpr 和普通 if 的差别}

普通 \code{if} 的两个分支都必须能通过类型检查：

\begin{lstlisting}[language=C++,caption={普通 if 不能隐藏类型错误}]
if (condition) {
    value.size();
} else {
    value + 1;
}
\end{lstlisting}

哪怕运行时不会走某个分支，编译器也要确认它语法成立。\code{if constexpr} 则会在编译期选择分支，未选择分支不实例化：

\begin{lstlisting}[language=C++,caption={if constexpr 按类型选择}]
template <typename T>
std::string describe_value(const T& value)
{
    if constexpr (requires { value.size(); }) {
        return "size=" + std::to_string(value.size());
    } else {
        return "scalar";
    }
}
\end{lstlisting}

这就是它适合泛型代码的原因：不同类型可以走不同代码路径，而不适用的路径不会强行检查。

\section{type traits 如何被 concept 包装}

旧代码常写 traits：

\begin{lstlisting}[language=C++,caption={traits 版本}]
template <typename T>
constexpr bool IS_SMALL_TRIVIAL_VALUE =
    std::is_trivially_copyable_v<T> && sizeof(T) <= 16;
\end{lstlisting}

C++20 可以把它包装成 concept：

\begin{lstlisting}[language=C++,caption={concept 包装类型条件}]
template <typename T>
concept SmallTrivialValue =
    std::is_trivially_copyable_v<T> && sizeof(T) <= 16;
\end{lstlisting}

之后接口可以直接写：

\begin{lstlisting}[language=C++,caption={接口用 concept 表达意图}]
template <SmallTrivialValue T>
void pass_by_value(T value);
\end{lstlisting}

traits 像螺丝和齿轮，concept 像对外写清楚的门牌。业务代码优先读门牌，底层库才经常直接摆弄齿轮。

\section{FixedString 案例}

一个固定字符串类型可以保存长度进入类型：

\begin{lstlisting}[language=C++,caption={固定字符串雏形}]
template <std::size_t N>
struct FixedString {
    std::array<char, N> chars{};

    constexpr explicit FixedString(const char (&text)[N])
    {
        std::copy_n(text, N, this->chars.begin());
    }

    [[nodiscard]] constexpr bool empty() const noexcept
    {
        return N <= 1;
    }
};
\end{lstlisting}

如果写 \code{FixedString{"while"}}，长度 \code{N} 来自字符串字面量，编译器能知道它。这个例子说明哪些信息能进编译期：字面量长度、模板参数、类型属性。用户运行时输入的一行文本，仍然不能变成模板参数，除非你重新编译程序。

\section{编译期不是质量保证的全部}

把错误提前到编译期很好，但不能代替运行期防御。端口范围函数可以 \code{static_assert(is_valid_port(80))}，但用户配置里的 \code{70000} 仍然要运行期报错。编译期工具负责固定世界，运行期错误处理负责外部世界。两本账都要有。
